\begin{figure}[t]
\centering
\begin{tikzpicture}[font=\small]

% Trakt som smalnar: fire-kant som er brei oppe og smal nede
\coordinate (tl) at (-32mm, 0);
\coordinate (tr) at ( 32mm, 0);
\coordinate (br) at ( 9mm, -48mm);
\coordinate (bl) at (-9mm, -48mm);

\fill[black!8] (tl) -- (tr) -- (br) -- (bl) -- cycle;
\draw[thick] (tl) -- (tr) -- (br) -- (bl) -- cycle;

% Tre nivå med vassrette skiljelinjer
\draw[gray!60] (-21mm,-16mm) -- (21mm,-16mm);
\draw[gray!60] (-15mm,-32mm) -- (15mm,-32mm);

% Etikettar for dei tre stega
\node[font=\bfseries] at (0,-8mm) {Mysterium};
\node[font=\bfseries] at (0,-24mm) {Heuristikk};
\node[font=\bfseries] at (0,-40mm) {Algoritme};

% Pil som viser rørsla nedover (innsnevringa)
\draw[-{Latex[length=2mm]}, thick] (40mm,-2mm) -- (40mm,-46mm)
  node[midway, right, align=left, font=\footnotesize]{stadig\\snævrare\\og\\billegare};

\end{tikzpicture}
\caption{\textsc{Kunnskapstrakta.}
Roger Martins «knowledge funnel»: arbeidet rører seg frå det opne mysteriet,
gjennom ein heuristikk som gjev retning, til ein algoritme som er billeg å
køyre i stor skala. Design thinking selde sjølve trakta som metode --- forma,
rørsla, dei tre orda --- utan å levere den kunnskapen som skal fylle henne.
Trakta utan innhald flyttar ingenting nedover.}
\label{fig:knowledge-funnel}
\end{figure}
